\documentclass[11pt,a4paper]{article}
\usepackage[utf8]{inputenc}
\usepackage[T1]{fontenc}
\usepackage{amsmath,amssymb}
\usepackage{geometry}
\geometry{top=2.5cm, bottom=2.5cm, left=2.5cm, right=2.5cm}
\usepackage{parskip}

\begin{document}

{\noindent\large\textbf{GMGD200}}

\vspace{0.2cm}
{\noindent\Large\textbf{Øving 1 - Koordinatsystemer}}

\vspace{0.6cm}
{\noindent\large\textbf{Sfæriske koordinater}}

\vspace{0.3cm}
\noindent Gitt posisjon for punkt $P$:
\begin{align*}
\phi_P &= 63^\circ 25' 10.50000'' \\
\Delta\lambda &= -0^\circ 19' 36.60000'' \\
h_P &= 135.500\text{ m}
\end{align*}

\begin{enumerate}
    \item 
    \begin{enumerate}
        \item $\lambda_P$ er gitt relativt akse III i det norske NGO1948 systemet. Sentralmeridianen $\lambda_0$ for akse III er gitt: $\lambda_0 = 10^\circ 43' 22.5''$
        \begin{enumerate}
            \item[i.] Beregn $\lambda_P = \lambda_0 + \Delta\lambda$ relativt Greenwich meridianen.
            \item[ii.] Uttrykk $\phi$ og $\lambda$ i grader og radianer.
            \item[iii.] Hvor mange grader er 1 radian?
            \item[iv.] Hvor mange radianer er $90^\circ$?
            \item[v.] Velg $R = 6378\,137\text{ m}$. Betrakt $\{\phi_P, \lambda_P, h_P\}$ som sfæriske koordinater og beregn de jordsentriske (ECEF) koordinatene $\{X_P, Y_P, Z_P\}$ for punkt $P$.
            \item[vi.] Hvilken enhet har de jordsentriske (ECEF) koordinatene?
        \end{enumerate}
        \item Legg til 100 meter på $X$-koordinaten til punkt $P$. Gi det nye punktet navnet $P'$ med tilhørende koordinat $X'$. Beregn fra $\{X', Y, Z\}$ de sfæriske koordinatene $\phi, \lambda$ og $h$ til punktet $P'$.
        \item Gjenta prosedyren to ganger ved å legge til 100 meter på henholdsvis $Y$- og $Z$-koordinaten $\{X, Y', Z\}$ og $\{X, Y, Z'\}$.
        \item Hvor stor er differansen $(P - P')$ uttrykt i sfæriske koordinater når $P'$ er henholdsvis $\{X', Y, Z\}$, $\{X, Y', Z\}$ og $\{X, Y, Z'\}$?
    \end{enumerate}
\end{enumerate}

\newpage

{\noindent\large\textbf{Ellipsoidiske koordinater}}

\vspace{0.3cm}
\begin{enumerate}
    \item[2.] 
    \begin{enumerate}
        \item Datumet EUREF89 benytter en GRS80 ellipsoide med store halvakse $a = 6378\,137\text{ m}$ og eksentrisitet $e^2 = 6694380.02290 \cdot 10^{-9}$. Beregn følgende ellipsoidiske størrelser:
        \begin{enumerate}
            \item[i.] Lille halvakse ($b$)
            \item[ii.] Flattrykning ($f$)
            \item[iii.] Andre nummeriske eksentrisitet ($e'^2$)
            \item[iv.] Differansen mellom halvaksene ($a - b$)
        \end{enumerate}
        \item Vis at den sfæriske koordinattransformasjonen $\{\phi, \lambda, h\} \to \{X, Y, Z\}$ er et spesialtilfelle av ellipsoidiske koordinattransformasjonen. (hint: sett $a = b = R$).
        \item Betrakt koordinatene $\{\phi_P, \lambda_P, h_P\}$ til punkt $P$ som ellipsoidiske koordinater på en GRS80 ellipsoide. Beregn de tilhørende jordsentriske (ECEF) koordinatene $\{X_P, Y_P, Z_P\}$.
        \item Hvor stor er differansen mellom de ellipsoidiske koordinatene $\{X_P, Y_P, Z_P\}$ til $P$ og de sfæriske koordinatene $\{X_P, Y_P, Z_P\}$ til $P$ fra forrige oppgave?
        \item Vis at en endring i vinkel på $1 \times 10^{-9}\text{ rad}$ i $\phi$ utgjør en avstand på $6\text{ mm}$ på jordoverflaten.
        \item Beregn de ellipsoidiske koordinatene $\{\phi_P, \lambda_P, h_P\}$ ved iterasjon.
        \item Beregn de ellipsoidiske koordinatene $\{\phi_P, \lambda_P, h_P\}$ ved bruk av Bowrings metode.
        \item Hva er forskjellen mellom disse to metodene?
    \end{enumerate}
\end{enumerate}

\end{document}
